\documentclass{beamer}

% Theme selection
\usetheme{metropolis}
 % You can choose from various themes like Berlin, AnnArbor, etc.

% Packages
\usepackage[utf8]{inputenc}
\usepackage{graphicx}
\usepackage{listings}
\usepackage{hyperref}
\usepackage{xcolor}

\setbeamertemplate{section page}{}
\setbeamertemplate{subsection page}{}

% Code Style
\lstset{
    language=C++,                % Set language to C++
    basicstyle=\ttfamily\footnotesize, % Monospace font, small size
    frame=shadowbox,             % Add a shadowed frame around the code
    rulesepcolor=\color{gray},   % Frame border color
    numbers=left,                % Add line numbers on the left
    numberstyle=\tiny\color{gray}, % Line number style
    breaklines=true,             % Automatically break long lines
    keywordstyle=\color{blue},   % Color for keywords
    commentstyle=\color{green!50!black}, % Color for comments
    stringstyle=\color{red},     % Color for strings
    showstringspaces=false,      % Do not show underscores in strings
}

% Title Page Details
\title{Facade Design Pattern}
\subtitle{Simplifying Complex Systems}
\author{  
Md Mehedi Hasan Maruf-2003037  \and
Md Shahidul Islam Dipu-2003040 \and \\
Eazdan Mostafa Rafin-2003055
}
\institute{Department of Computer Science and Engineering, RUET}
\date{\today}

% Begin Document
\begin{document}

% Title Slide
\begin{frame}
    \titlepage
\end{frame}

% Table of Contents
% \begin{frame}{Agenda}
%     \tableofcontents
% \end{frame}

% Section 1: Introduction

\begin{frame}{What is a Design Pattern?}
    \begin{itemize}
        \item Reusable solutions to common software design problems
        \item Not code, but templates for solving problems
        \item Enhance code maintainability and scalability
    \end{itemize}
    \vspace{0.5cm}
    \textbf{Types of Design Patterns:}
    \begin{itemize}
        \item \textbf{Creational Patterns}
        \item \textbf{Structural Patterns} 
        \item \textbf{Behavioral Patterns} 
    \end{itemize}
\end{frame}

% Section 2: Facade Design Pattern Overview

\begin{frame}{Facade Design Pattern}
    \begin{itemize}
        \item \textbf{Type:} Structural Design Pattern
        \item \textbf{Objective:}
        \begin{itemize}
            \item Provide a simplified interface to complex subsystems.
            \item Reduce dependencies between the client and subsystem classes.
        \end{itemize}
    \end{itemize}
\end{frame}


% Section 3: When to Use

\begin{frame}{When to Use Facade}
    \begin{itemize}
        \item Complex Subsystems with many interdependent classes
        \item Need to provide a simple interface for clients
        \item Layered systems requiring interaction with lower layers
        \item Integrating with legacy or third-party code
    \end{itemize}
\end{frame}

% Section 4: Structure

% \begin{frame}{UML Class Diagram}
%     \begin{center}
%         \includegraphics[width=0.8\linewidth]{uml_facade.png} % Ensure you have this image in your directory
%     \end{center}
%     \footnotesize{Figure: UML Class Diagram of Facade Pattern}
% \end{frame}

% Section 5: Components

\begin{frame}{Key Components}
    \begin{enumerate}
        \item \textbf{Facade:} Simplified interface for the client
        \item \textbf{Subsystem Classes:} Complex classes handling specific functionalities
        \item \textbf{Client:} Interacts only with the Facade
    \end{enumerate}
\end{frame}



% Code Example Slide
\begin{frame}[fragile]{C++ Code Example}
    \begin{lstlisting}
    // Facade Design Pattern in C++
    #include <iostream>
    using namespace std;
    
    class Light {
    public:
        void dim() { cout << "Lights dimmed.\n"; }
    };
    
    class MusicSystem {
    public:
        void play() { cout << "Playing music.\n"; }
    };
    
    class HomeTheaterFacade {
        Light light;
        MusicSystem music;
    public:
        void activateMovieMode() {
            cout << "Activating Movie Mode...\n";
            light.dim();
            music.play();
        }
    };
    
    int main() {
        HomeTheaterFacade theater;
        theater.activateMovieMode();
        return 0;
    }
    \end{lstlisting}
    \end{frame}
% Section 6: Benefits and Drawbacks

\begin{frame}{Pros}
    \begin{itemize}
        \item Simplifies interactions with complex systems
        \item Promotes loose coupling between client and subsystem
        \item Enhances code readability and maintainability
        \item Facilitates easier testing and debugging
    \end{itemize}
\end{frame}


\begin{frame}{Cons}
    \begin{itemize}
        \item \textbf{Hides Too Much:} Can block access to advanced features.
        \item \textbf{Overdependence:} Makes it harder to work without the Facade.
        \item \textbf{Maintenance Overhead:} Needs regular updates to stay in sync.
    \end{itemize}
\end{frame}



% Section 12: Real-World Applications

\begin{frame}{Real-World Applications}
    \begin{itemize}
        \item \textbf{Home Theater Systems:} One button dims lights, starts the projector, and plays your movie.
        \item \textbf{Computer Start-Up:} Turns chaos into a smooth boot-up process.
        \item \textbf{Online Shopping:} Handles payments, inventory, and shipping in the background.
    \end{itemize}
\end{frame}

\begin{frame}{Summary}
    \centering
    The Facade Design Pattern simplifies complex systems by providing a clean and easy-to-use interface, making your code more readable, maintainable, and less dependent on the intricate details of subsystems.
\end{frame}


\end{document}
